\section{子空间的有理性}

记号约定:

\begin{itemize}
	\item $K$是除环,$K^{\prime}$是其子除环,$\rho:K^{\prime}\hookrightarrow K$是典范嵌入.
	\item $V$是右$K$-向量空间,$V^{\prime}$是$V$的$K^{\prime}$-结构.
\end{itemize}

\begin{definition}{有理向量}{}
	设$x\in V$.若$x\in V^{\prime}$,则称$x$在$K^{\prime}$上有理.
\end{definition}

\begin{definition}{有理子空间}{}
	设$W$是$V$的子空间.记$W^{\prime}=\rho_{\ast}\left(W\right)\cap V$.若典范映射$\rho_{\ast}\left(W\right)\to W$是$K$-线性同构,则称$W$是$V$的$K^{\prime}$-有理子空间.
\end{definition}

\begin{proposition}{有理向量的坐标判别法}{}
	设$\left(v_{i}^{\prime}\right)_{i\in I}$是$V^{\prime}$的基,$x=\sum\limits_{i\in I}v_{i}^{\prime}\xi_{i}\in V$,则$x$在$K^{\prime}$上有理$\iff$ $\left(\xi_{i}\right)_{i\in I}$是$\rho\left(K^{\prime}\right)$的元素族.
\end{proposition}

\begin{proposition}{诱导$K^{\prime}$-结构}{}
	设$W$是$V$的子空间.若$W$在$K^{\prime}$上有理,则$\rho_{\ast}\left(W\right)\cap V^{\prime}$是$W$的$K^{\prime}$-结构.
\end{proposition}

\begin{definition}{诱导的$K^{\prime}$-结构}{}
	设$W$是$V$的在$K^{\prime}$上有理的子空间.我们称$\rho_{\ast}\left(W\right)\cap V^{\prime}$是$V^{\prime}$在$W$上诱导的$K^{\prime}$-结构.
\end{definition}

\begin{proposition}{$K^{\prime}$-子空间和$K^{\prime}$-有理子空间的对应}{}
	设$\mathcal{S}\left(V^{\prime}\right)$是$V^{\prime}$的子空间组成的集合,$\mathcal{S}_{\text{rat}}\left(V\right)$是$V$的在$K^{\prime}$上有理的子空间组成的集合,则
	\begin{align*}
		\phi:\mathcal{S}\left(V^{\prime}\right)\to\mathcal{S}_{\text{rat}}\left(V\right),W\mapsto\langle W\rangle
	\end{align*}
	是双射,其逆为$\psi:\mathcal{S}_{\text{rat}}\left(V\right)\to\mathcal{S}\left(V^{\prime}\right),W\mapsto \rho_{\ast}\left(W\right)\cap V^{\prime}$.
\end{proposition}

\begin{corollary}{}{}
	设$\left(W_{i}\right)_{i\in I}$是$V$的一族在$K^{\prime}$上有理的子空间,则$\bigcap\limits_{i\in I}W_{i}$和$\sum\limits_{i\in I}W_{i}$在$K^{\prime}$上有理.
\end{corollary}

\begin{definition}{$K^{\prime}$-有理基集}{}
	设$B$是$V$的基.若$B$的每个元素都在$K^{\prime}$上有理,则称$B$是$K^{\prime}$-有理基集.
\end{definition}

\begin{corollary}{}{}
	若$B$是$V$的$K^{\prime}$-基集,则$B$是$V^{\prime}$的基.
\end{corollary}